\section*{Заключение}
\addcontentsline{toc}{section}{Заключение}

В ходе выполнения курсовой работы было спроектировано и реализовано
веб-приложение для управления бронированием столов и заказами
кафе~«Table~\&~Time». Приложение представляет собой полноценную
клиент-серверную систему с разграничением прав доступа, интерактивным
интерфейсом и автоматизированной бизнес-логикой.

В результате работы достигнуты все поставленные цели:

\begin{itemize}
  \item разработана и задокументирована реляционная модель данных,
        включающая восемь сущностей и семь перечислений, отражающих
        предметную область заведения общественного питания;
  \item реализован REST API на основе фреймворка NestJS с системой
        аутентификации и авторизации на JWT-токенах, передаваемых через
        HTTP-only куки, и разграничением прав для четырёх ролей
        пользователей;
  \item разработан интерактивный клиент на Next.js~15 с использованием
        App Router, включающий карту столов с цветовой индикацией
        статусов в реальном времени, форму бронирования с проверкой
        доступности, страницы меню, оформления и отслеживания заказов,
        личный кабинет и полнофункциональную административную панель;
  \item реализована система автоматического управления статусами столов
        на основе таймеров в памяти процесса с восстановлением состояния
        после перезапуска сервера;
  \item приложение контейнеризовано с помощью Docker и развёрнуто
        в облачной инфраструктуре: серверная часть~--- на платформе
        Railway, клиентская~--- на Vercel; организован проксирующий
        слой через механизм rewrites Next.js для обхода браузерных
        ограничений на межсайтовые cookie.
\end{itemize}

В процессе разработки были применены современные подходы и инструменты:
декларативное управление схемой базы данных через Prisma Migrate,
типобезопасная разработка на TypeScript, декомпозиция серверной логики
на модули по принципам NestJS, клиентская маршрутизация с защитой
страниц на уровне компонентов, а также непрерывная доставка (CD) через
интеграцию GitHub с облачными платформами.

Разработанное приложение полностью соответствует сформулированным
требованиям и может быть использовано в реальных условиях заведения
общественного питания. Полученный опыт проектирования и реализации
полностековых веб-приложений, включая вопросы развёртывания,
безопасности и управления сессиями, представляет практическую ценность
для дальнейшей профессиональной деятельности.